%% Solutions for Chapter 12: Decidability

\chapter{Decidability --- Solutions}

\begin{enumerate}[{12.}1. ]

%%------------------------------------------------------------
\item \textbf{Verify assertions in proof of Theorem 12.1 concerning Theorem 2.7.}

The point used in Theorem 12.1 is this: if $A$ accepts any string at all, then
$A$ accepts one of length $< n$, where $n=\|S\|$.

Let
\[
A=\langle \Sigma,S,s_0,\delta,F\rangle
\]
and suppose $x\in L(A)$ with $|x|\geq n$. By Theorem 2.7, $L(A)$ is the union
of classes of a finite-rank right congruence; in particular, we may use
$R^A$, whose rank is at most $\|S\|=n$. Write
\[
x=a_1a_2\cdots a_m
\]
and consider the $n+1$ prefixes
\[
\lambda,\ a_1,\ a_1a_2,\ \ldots,\ a_1a_2\cdots a_n.
\]
Since $R^A$ has at most $n$ classes, two of these prefixes must be related by
$R^A$. So we can write
\[
x=uvw
\]
with $v\neq\lambda$ and $u\ R^A\ uv$. Because $R^A$ is a right congruence,
\[
uw\ R^A\ uvw=x.
\]
Now $L(A)$ is a union of $R^A$-classes, and $x\in L(A)$, so $uw\in L(A)$ as
well. Thus we have found a shorter accepted string. Repeating this shortening
process eventually produces an accepted string of length $<n$.

That is exactly the assertion needed in the proof of Theorem 12.1. \qed

%%------------------------------------------------------------
\item \textbf{Prove Lemma 12.1.}

Let
\[
M=\{x\mid x\in L(A)\wedge |x|\geq n\}.
\]
If $M\neq\emptyset$, then the set
\[
\{|x|\mid x\in M\}\subseteq \NN
\]
is nonempty, so by the well-ordering principle it has a least element. Hence
$M$ contains a string of minimal length; call it $x_m$.

It remains to prove that $|x_m|<2n$. Suppose, for contradiction, that
$|x_m|\geq 2n$. Since $x_m\in L(A)$ and $|x_m|\geq n$, the pumping lemma
(Theorem 2.3) gives a factorization
\[
x_m=uvw
\]
with
\[
|uv|\leq n,\qquad |v|>0,
\]
and
\[
uv^iw\in L(A)\qquad\text{for all }i\in\NN.
\]
Take $i=0$. Then
\[
uw\in L(A).
\]
Also,
\[
|uw|=|x_m|-|v|\geq |x_m|-n\geq n,
\]
since $|v|\leq |uv|\leq n$. Thus $uw\in M$. But $|v|>0$, so
\[
|uw|<|x_m|,
\]
contradicting the minimality of $x_m$.

Therefore $|x_m|<2n$. \qed

%%------------------------------------------------------------
\item \textbf{Number of nonfinal states in minimal DFA $\leq$ nonfinal states in any machine $B$ with $L(B)=L$.}

\textbf{Proof}: Let $M$ be the minimal DFA for $L$. By Nerode's theorem, the
states of $M$ correspond exactly to the classes of $R_L$. In particular, the
nonfinal states of $M$ correspond exactly to the classes $C$ for which every
string in $C$ is \emph{not} in $L$.

For each such nonfinal class $C$, choose a representative string $x_C\notin L$.
Now let
\[
B=\langle \Sigma,S_B,s_B,\delta_B,F_B\rangle
\]
be any DFA with $L(B)=L$. Since $x_C\notin L$, the state
\[
\DBAR_B(s_B,x_C)
\]
must be nonfinal.

We claim that different nonfinal classes of $R_L$ must lead to different states
of $B$. Suppose $C_1\neq C_2$ but
\[
\DBAR_B(s_B,x_{C_1})=\DBAR_B(s_B,x_{C_2}).
\]
Since $C_1$ and $C_2$ are distinct Nerode classes, there exists a string $z$
such that exactly one of $x_{C_1}z$ and $x_{C_2}z$ belongs to $L$. But starting
from the same state of $B$, reading the same suffix $z$ must give the same
accept/reject answer for both strings, a contradiction.

Thus distinct nonfinal Nerode classes determine distinct nonfinal states of
$B$. Therefore the number of nonfinal states in $M$ is less than or equal to the
number of nonfinal states in $B$. \qed

%%------------------------------------------------------------
\item \textbf{Decidability of $L(A_1)\subseteq L(A_2)$.}

\textbf{Algorithm}:
\begin{enumerate}
\item Construct DFA $A_2^c$ for $\overline{L(A_2)}$ (complement: swap final/nonfinal states of $A_2$'s complete DFA).
\item Construct DFA $A^\cap$ for $L(A_1)\cap\overline{L(A_2)}$ (product machine).
\item $L(A_1)\subseteq L(A_2)$ iff $L(A_1)\cap\overline{L(A_2)}=\emptyset$ iff $L(A^\cap)=\emptyset$ (decidable by Theorem 12.1). \qed
\end{enumerate}

%%------------------------------------------------------------
\item \textbf{Decidability of $L(A)$ is cofinite.}

$L(A)$ is cofinite iff $\overline{L(A)}$ is finite.

\textbf{Algorithm}:
\begin{enumerate}
\item Construct DFA $A^c$ for $\overline{L(A)}$ (complement).
\item $\overline{L(A)}$ is finite iff $L(A^c)$ is finite iff the DFA $A^c$ (after removing unreachable states) has no cycle reachable from the start state that also leads to a final state (Theorem 12.2 analog).
\item Check for cycles in $A^c$ among useful states (BFS/DFS). Decidable. \qed
\end{enumerate}

%%------------------------------------------------------------
\item \textbf{Decidability: does $L(A)$ contain any string of length $> 1228$?}

\textbf{Algorithm}: Build the DFA
\[
B=\langle \Sigma,\{r_0,r_1,\ldots,r_{1229}\},r_0,\delta_B,\{r_{1229}\}\rangle,
\]
where for every $a\in\Sigma$,
\[
\delta_B(r_i,a)=r_{\min\{i+1,1229\}}.
\]
Then
\[
L(B)=\{x\in\Sigma^*\mid |x|\geq 1229\},
\]
that is, the strings of length greater than $1228$.

Now form the product DFA $A^\cap$ for
\[
L(A)\cap L(B).
\]
Then
\[
L(A)\text{ contains a string of length }>1228
\]
iff
\[
L(A^\cap)\neq\emptyset.
\]
By Theorem 12.1, emptiness for DFAs is decidable. Hence this question is
decidable. \qed

%%------------------------------------------------------------
\item \textbf{Decidability: does $A$ accept any even-length strings?}

\textbf{Algorithm}:
\begin{enumerate}
\item Construct DFA $A_{even}$ that recognizes $\{x \mid |x|$ is even$\}$ (a simple 2-state DFA alternating between even/odd states).
\item Construct product DFA $A^\cap$ for $L(A)\cap L(A_{even})$.
\item $A$ accepts an even-length string iff $L(A^\cap)\neq\emptyset$. Decidable by Theorem 12.1. \qed
\end{enumerate}

%%------------------------------------------------------------
\item \textbf{Decidability: $R_1$ and $R_2$ represent complementary languages.}

\textbf{Algorithm}:
\begin{enumerate}
\item Convert regular expressions $R_1$ and $R_2$ to DFAs $A_1$ and $A_2$.
\item Check $L(A_1) = \overline{L(A_2)}$: equivalently, check $L(A_1)\subseteq\overline{L(A_2)}$ and $\overline{L(A_2)}\subseteq L(A_1)$.
\item By Exercise 12.4, each subset relation is decidable. Equivalently, check $L(A_1) \cup L(A_2) = \Sigma^*$ and $L(A_1)\cap L(A_2)=\emptyset$.
\item Both conditions are decidable (using complement, intersection, and emptiness). \qed
\end{enumerate}

%%------------------------------------------------------------
\item \textbf{Decidability: $R_1$ and $R_2$ describe common strings.}

\textbf{Algorithm}: $L(R_1)\cap L(R_2)\neq\emptyset$?
\begin{enumerate}
\item Convert $R_1,R_2$ to DFAs $A_1,A_2$.
\item Compute the product DFA $A^\cap$ for $L(A_1)\cap L(A_2)$.
\item Check $L(A^\cap)\neq\emptyset$ (decidable by Theorem 12.1). \qed
\end{enumerate}

%%------------------------------------------------------------
\item \textbf{Decidability: is there a DFA with $< 31$ states accepting $L(R_1)$?}

\textbf{Algorithm}:
\begin{enumerate}
\item Convert $R_1$ to a DFA $A$ (via NFA construction).
\item Minimize $A$ to get $A_{min}$.
\item There is a DFA with $<31$ states accepting $L(R_1)$ iff $|A_{min}| < 31$.
\item Count the states of $A_{min}$: decidable. \qed
\end{enumerate}

%%------------------------------------------------------------
\item \textbf{Decidability: is there a DFA with $>31$ states accepting $L(R_1)$?}

\textbf{One-step algorithm}: Yes, always. If $A_{min}$ has $k$ states, then we can construct a DFA with $k + \ell$ states for any $\ell > 0$ by adding dead (sink) states. If $k \leq 31$, add $32-k$ unreachable dead states. If $k > 31$, $A_{min}$ itself already has $>31$ states.

Therefore, for \emph{any} regular language, there always exists a DFA with $>31$ states accepting it. The answer is \textbf{always yes}. \qed

%%------------------------------------------------------------
\item \textbf{Decidability: $\exists$ $\lambda$-NDFA with $\leq 1$ final state accepting $R$.}

\textbf{Algorithm}:
\begin{enumerate}
\item Convert $R$ to a $\lambda$-NDFA $N$ using the Kleene/Thompson construction. This construction produces NDFAs with exactly one final state. So the answer is always yes for any regular expression $R$. \qed
\end{enumerate}

%%------------------------------------------------------------
\item \textbf{Decidability: $\exists$ NDFA (without $\lambda$-moves) with $\leq 1$ final state accepting $L(A)$.}

Let $L=L(A)$. We will show that such an NDFA exists iff either
\[
\lambda\notin L,
\]
or else
\[
\lambda\in L \quad\text{and}\quad L\cdot L\subseteq L.
\]

\textbf{Necessity}: Suppose
\[
N=\langle \Sigma,S,s_0,\delta,\{f\}\rangle
\]
is an NDFA without $\lambda$-moves that accepts $L$ and has at most one final
state.

If $\lambda\notin L$, there is nothing more to prove. So assume $\lambda\in L$.
Then the start state itself must be final, and since there is at most one final
state we have
\[
f=s_0.
\]
Therefore every accepted word labels a path from $s_0$ back to $s_0$. If
$x,y\in L$, we may follow a path for $x$ from $s_0$ back to $s_0$ and then a
path for $y$ from $s_0$ back to $s_0$. Hence
\[
xy\in L.
\]
So whenever $\lambda\in L$, the language must be closed under concatenation.

\textbf{Sufficiency when $\lambda\notin L$}: Since regular languages are closed
under reversal, construct a DFA for $L^r$ and reverse all its arrows. This gives
an NDFA $M$ for $L$ that has one final state $q$ and possibly several start
states collected in a set $I$. Because $\lambda\notin L$, the state $q$ is not a
member of $I$.

Now replace the set $I$ by a single new start state $s$, and for each
$a\in\Sigma$ define
\[
\delta(s,a)=\bigcup_{i\in I}\delta_M(i,a).
\]
Leave all other transitions unchanged and keep $q$ as the unique final state.
Every accepting path in the old machine is simulated by first choosing the same
first step from $s$, so the new machine accepts exactly $L$. Thus the answer is
Yes whenever $\lambda\notin L$.

\textbf{Sufficiency when $\lambda\in L$ and $L$ is closed under concatenation}:
Start again with the reversed machine $M$ above. Now $q\in I$. Make $q$ the
only start state, and redefine its outgoing transitions by
\[
\delta'(q,a)=\bigcup_{i\in I}\delta_M(i,a).
\]
All other transitions remain as in $M$, and $q$ is still the unique final state.

Every word of $L$ is accepted by this new machine, because its first move can
imitate the first move from whichever old start state in $I$ was used before.
Conversely, any accepting path in the new machine starts at $q$ and ends at $q$.
Whenever the path returns to $q$, split it there. Each resulting segment is a
word that was already accepted by the old multi-start machine, hence is in $L$.
Since $L$ is closed under concatenation, the whole input word also belongs to
$L$. So this new machine again accepts exactly $L$ and has one final state.

\textbf{Decision procedure}:
\begin{enumerate}
\item Check whether $\lambda\in L(A)$ by seeing whether the start state of $A$ is
final.
\item If $\lambda\notin L(A)$, answer Yes.
\item If $\lambda\in L(A)$, construct an automaton for $L(A)\cdot L(A)$ and use
Exercise 12.4 to decide whether
\[
L(A)\cdot L(A)\subseteq L(A).
\]
\end{enumerate}
This inclusion is exactly the closure-under-concatenation condition above, so
the question is decidable. \qed

%%------------------------------------------------------------
\item \textbf{Decidability of $R_1 = R_2$ (same language).}

\textbf{Algorithm}:
\begin{enumerate}
\item Convert $R_1, R_2$ to DFAs $A_1, A_2$.
\item Minimize both: $A_1^{min}, A_2^{min}$.
\item $L(R_1)=L(R_2)$ iff $A_1^{min} \cong A_2^{min}$ (isomorphic as DFAs). Check isomorphism: decidable in polynomial time. \qed
\end{enumerate}

%%------------------------------------------------------------
\item \textbf{Decidability: $R_1$ and $R_2$ generate the same right congruence ($R_{L_1}=R_{L_2}$).}

\textbf{Algorithm}: Convert $R_1$ and $R_2$ to minimal DFAs
\[
M_1=\langle\Sigma,S_1,s_1,\delta_1,F_1\rangle,\qquad
M_2=\langle\Sigma,S_2,s_2,\delta_2,F_2\rangle.
\]
By Exercise 3.40, since these machines are minimal,
\[
R_{L_1}=R^{M_1},\qquad R_{L_2}=R^{M_2}.
\]
Therefore
\[
R_{L_1}=R_{L_2}
\]
iff
\[
R^{M_1}=R^{M_2}.
\]
Because $M_1$ and $M_2$ are connected, this happens exactly when there is a
bijection
\[
\mu:S_1\to S_2
\]
such that
\[
\mu(s_1)=s_2
\]
and
\[
\mu(\delta_1(s,a))=\delta_2(\mu(s),a)
\]
for every state $s$ and letter $a$. Notice that the final states are
\emph{irrelevant}; we are comparing the right congruences, not the languages.

So the decision procedure is: minimize both machines, forget which states are
final, and test whether the two rooted transition graphs are isomorphic. This
is a finite check, and hence decidable. \qed

%%------------------------------------------------------------
\item \textbf{Prove Theorem 12.5.}

Theorem 12.5 states that it is decidable whether two right-linear grammars
\[
G_1=\langle\Omega_1,\Sigma,S_1,P_1\rangle,\qquad
G_2=\langle\Omega_2,\Sigma,S_2,P_2\rangle
\]
are equivalent.

Use Lemma 8.2 to build NDFAs $A_{G_1}$ and $A_{G_2}$ for the two grammars.
Convert these NDFAs to DFAs $A_1$ and $A_2$. Then
\[
L(G_1)=L(G_2)
\]
iff
\[
L(A_1)=L(A_2).
\]
By Theorem 12.3, equivalence of DFAs is decidable. Therefore equivalence of the
two right-linear grammars is decidable as well. \qed

%%------------------------------------------------------------
\item \textbf{Efficient algorithm for $\{\DBAR(s_0,y)\mid y\in\Sigma^*,\ n\leq|y|<2n\}$.}

\textbf{Algorithm}: Compute the sets
\[
T_i=\{\DBAR(s_0,y)\mid y\in\Sigma^i\}
\]
recursively.
\begin{enumerate}
\item Set $T_0=\{s_0\}$.
\item For $i=0,1,\ldots,2n-2$, compute
\[
T_{i+1}=\{\delta(s,a)\mid s\in T_i,\ a\in\Sigma\}.
\]
\item Return
\[
T_n\cup T_{n+1}\cup\cdots\cup T_{2n-1}.
\]
\end{enumerate}
Each $T_i$ is exactly the set of states reached by strings of length $i$, so the
returned union is the desired set. The procedure halts because it performs only
$2n-1$ finite computations. \qed

%%------------------------------------------------------------
\item \textbf{Would intersecting $\{\DBAR(s_0,y)\mid 5n\leq|y|<6n\}$ with $F$ always work?}

\textbf{Answer}: Yes, it would always produce the correct answer.

\textbf{Justification}: If
\[
\{\DBAR(s_0,y)\mid 5n\leq|y|<6n\}\cap F\neq\emptyset,
\]
then some accepted string has length at least $5n$, hence at least $n$, and the
argument used in Theorem 12.2 shows immediately that $L(A)$ is infinite.

For the converse, assume $L(A)$ is infinite. Then Theorem 12.2 tells us that
there is an accepted string $x$ with
\[
n\leq |x|<2n.
\]
Apply the pumping lemma to $x$, writing
\[
x=uvw
\]
with
\[
|uv|\leq n,\qquad |v|>0,
\]
and
\[
uv^iw\in L(A)\qquad\text{for all }i\in\NN.
\]
Choose the least $i$ such that
\[
|uv^iw|\geq 5n.
\]
Since $|v|\leq n$, minimality of $i$ gives
\[
|uv^iw|<5n+n=6n.
\]
Thus some accepted string has length in the range
\[
5n\leq |y|<6n.
\]
So this alternative interval also works. \qed

%%------------------------------------------------------------
\item \textbf{Decidability of equivalence of two Mealy machines.}

\textbf{Algorithm}:
\begin{enumerate}
\item Build the product graph on pairs of states $(q_1,q_2)$.
\item Start from the pair of initial states.
\item For each reachable pair $(q_1,q_2)$ and each input symbol $a$, compare the
two output symbols produced on $a$.
\item If a mismatch is ever found, the machines are not equivalent.
\item Otherwise continue to the pair of next states reached on $a$.
\end{enumerate}
Because the product graph is finite, this search halts. It answers Yes exactly
when every reachable pair produces the same output on every input symbol, which
is precisely equivalence of the two Mealy machines. \qed

%%------------------------------------------------------------
\item \textbf{Decidability of equivalence of two Moore machines.}

\textbf{Algorithm}: Use the product graph on pairs of states.
\begin{enumerate}
\item Start from the pair of initial states and first compare their output
symbols.
\item For every reachable pair $(q_1,q_2)$, compare $\omega_1(q_1)$ and
$\omega_2(q_2)$.
\item If a mismatch is found, the machines are not equivalent.
\item Otherwise follow matching input symbols to successor pairs and continue.
\end{enumerate}
The product graph is finite, so this procedure halts. It is correct because
Moore outputs are determined by the states themselves. \qed

%%------------------------------------------------------------
\item \textbf{Decide whether $R$ represents strings of length $> 28$, without finite automata.}

\textbf{Algorithm}: Compute a capped maximum-length function
\[
m(R)\in\{-\infty,0,1,\ldots,29\},
\]
where $m(R)=29$ means ``$R$ represents some string of length $>28$.'' Define it
recursively:

\begin{enumerate}
\item $m(\emptyset)=-\infty$.
\item $m(\lambda)=0$.
\item $m(a)=1$ for $a\in\Sigma$.
\item $m(R_1\cup R_2)=\max\{m(R_1),m(R_2)\}$.
\item For concatenation:
\[
m(R_1R_2)=
\left\{
\begin{array}{ll}
-\infty & \text{if }m(R_1)=-\infty\text{ or }m(R_2)=-\infty,\\
29 & \text{if }m(R_1)=29\text{ or }m(R_2)=29\text{ or }m(R_1)+m(R_2)>28,\\
m(R_1)+m(R_2) & \text{otherwise.}
\end{array}
\right.
\]
\item For star:
\[
m(R_1^*)=
\left\{
\begin{array}{ll}
0 & \text{if }L(R_1)\subseteq\{\lambda\},\\
29 & \text{otherwise.}
\end{array}
\right.
\]
\end{enumerate}

Then $R$ represents some string of length $>28$ iff $m(R)=29$. Since this
recursion follows the finite parse tree of $R$, it always halts. \qed

%%------------------------------------------------------------
\item \textbf{Decide whether $L(G)$ contains a string of length $>28$ for right-linear grammar $G$.}

\textbf{Algorithm} (without finite automata): Regard the grammar as a weighted
directed graph. For each production of the form
\[
A\to xB
\]
place an edge $A\to B$ of weight $|x|$, and for each production
\[
A\to x
\]
place an edge from $A$ to a terminal sink $\bot$ of weight $|x|$.

First discard useless nonterminals by the method of Theorem 12.6. Now there are
two possibilities.

\begin{enumerate}
\item There is a reachable productive cycle whose total weight is positive.
Then the cycle can be traversed arbitrarily many times before exiting to
$\bot$, so $L(G)$ contains arbitrarily long strings, and in particular a string
of length $>28$.

\item No such cycle exists. Then, after collapsing zero-weight cycles, the
useful part of the graph is acyclic. The maximum length of a string derivable
from each nonterminal can then be computed by dynamic programming in reverse
topological order:
\[
\ell(A)=\max\{|x|,\ |x|+\ell(B)\},
\]
taking the maximum over productions $A\to x$ and $A\to xB$. The answer is Yes
iff $\ell(S)>28$.
\end{enumerate}

Every step is finite, so the question is decidable. \qed

%%------------------------------------------------------------
\item \textbf{Prove $Z_0\subseteq Z_1\subseteq\cdots\subseteq Z_n\subseteq\cdots\subseteq\Omega$.}

In the proof of Theorem 12.6 (decidability of $L(G)=\emptyset$ for CFG $G$), $Z_i$ is the set of nonterminals reachable from the start symbol $S$ in $i$ or fewer derivation steps.

$Z_0 = \{S\}$ (only the start symbol).
$Z_{i+1} = Z_i \cup \{B \mid \exists A\in Z_i, \exists \alpha: A\to\alpha\in P, B\in\alpha\}$ (all nonterminals appearing in right-hand sides of productions of nonterminals in $Z_i$).

\textbf{Proof}: $Z_0\subseteq Z_1$ since $Z_1 = Z_0\cup(\cdots) \supseteq Z_0$. By induction: if $Z_k\subseteq Z_{k+1}$, then $Z_{k+1}\subseteq Z_{k+2}$ since any nonterminal in $Z_{k+1}$ is reachable in $k+1$ steps, hence reachable in $k+2$ steps as well. The chain is bounded above by $\Omega$ (which is finite). \qed

%%------------------------------------------------------------
\item \textbf{If $Z_m = Z_{m+1}$, then $Z_m$ is the set of all nonterminals reachable from $S$.}

\textbf{Proof}: If $Z_m = Z_{m+1}$, then no new nonterminals are added in step $m+1$. This means: every nonterminal appearing in the right-hand side of any production of any nonterminal in $Z_m$ is already in $Z_m$.

\textbf{Claim}: $Z_m$ is closed under reachability, i.e., every nonterminal reachable from $S$ is in $Z_m$.

Suppose $B$ is reachable from $S$ (i.e., $S\DRightarrow \alpha B\beta$ for some $\alpha,\beta$). The shortest derivation has some length $k$. By monotonicity ($Z_0\subseteq Z_1\subseteq\cdots$), $B\in Z_k$. Since $Z_m=Z_{m+1}=Z_{m+2}=\cdots$ (the sequence stabilizes), $B\in Z_m$ for all $k\leq m$ already. If $k>m$, then by the stabilization, every nonterminal reachable in $k$ steps is already in $Z_m$. So $Z_m$ = all reachable nonterminals. \qed

%%------------------------------------------------------------
\item \textbf{Prove $\exists m\in\NN: Z_m = Z_{m+1}$.}

\textbf{Proof}: The sequence $Z_0\subseteq Z_1\subseteq\cdots\subseteq\Omega$ is a non-decreasing chain of subsets of $\Omega$. Since $|\Omega|$ is finite, the chain can increase at most $|\Omega|$ times. Therefore there exists $m\leq|\Omega|$ such that $Z_m=Z_{m+1}$ (the chain stabilizes). \qed

%%------------------------------------------------------------
\item \textbf{Definition of $W_i$ and monotonicity.}

\begin{enumerate}
\item In the proof of Theorem 12.6, $W_i$ is the set of nonterminals that can derive a terminal string in $i$ or fewer steps:
$W_0 = \emptyset$.
$W_{i+1} = W_i \cup \{A \mid \exists (A\to\alpha)\in P: \text{every symbol in }\alpha\in\Sigma\cup W_i\}$.
(A nonterminal $A$ can derive a terminal string in $\leq i+1$ steps if some production of $A$ has a right-hand side consisting only of terminals and nonterminals already known to derive terminal strings.)

\item \textbf{Proof of monotonicity}: $W_0=\emptyset\subseteq W_1$ (trivially). By induction: if $W_k\subseteq W_{k+1}$, then for any $A\in W_{k+1}$: $A$ derives a terminal string in $\leq k+1$ steps, hence in $\leq k+2$ steps as well, so $A\in W_{k+2}$. Thus $W_{k+1}\subseteq W_{k+2}$. The chain is bounded by $\Omega$. \qed
\end{enumerate}

%%------------------------------------------------------------
\item \textbf{If $W_m = W_{m+1}$, then $W_m$ is the set of all nonterminals producing terminal strings.}

\textbf{Proof}: If $A\DRightarrow^* w\in\Sigma^*$ in $k$ steps, then $A\in W_k$. Since $W_m=W_{m+k}$ for all $k\geq 0$ (by stabilization), $A\in W_m$. Conversely, every $A\in W_m$ derives a terminal string by definition. So $W_m$ is precisely the set of productive nonterminals. \qed

%%------------------------------------------------------------
\item \textbf{Prove $\exists m\in\NN: W_m = W_{m+1}$.}

\textbf{Proof}: By the definition of the sets $W_i$,
\[
W_0\subseteq W_1\subseteq W_2\subseteq\cdots\subseteq\Omega.
\]
Whenever $W_i\neq W_{i+1}$, at least one new nonterminal of the finite set
$\Omega$ has been added. Since $\Omega$ has only finitely many elements, this can
happen only finitely many times. Therefore for some $m$ no new symbol is added,
so $W_m=W_{m+1}$. \qed

%%------------------------------------------------------------
\item \textbf{Decidability questions about NDFA $A$ with $\lambda$-moves.}

\begin{enumerate}
\item \textbf{Exists string with no complete path}: Build the equivalent DFA
$A^d$. For any string $x$, the state of $A^d$ reached by $x$ is exactly the set
of states of $A$ at the ends of complete paths labeled $x$. Thus there is a
string with no complete path iff the empty subset is reachable in $A^d$.

\item \textbf{Exists string with a path to a nonfinal state}: In $A^d$, search
for a reachable subset that contains a nonfinal state of $A$. Such a subset
exists iff some string has at least one complete path ending in a nonfinal
state.

\item \textbf{Exists accepted string with all paths leading to final states}:
This asks whether there is a reachable subset $Q$ of $A^d$ such that
\[
Q\neq\emptyset,\qquad Q\subseteq F.
\]
Then $x$ is accepted (because $Q$ is nonempty and every element is final), and
every complete path on $x$ ends in a final state.

\item \textbf{All accepted strings have all complete paths to final states}:
This fails exactly when some reachable subset $Q$ of $A^d$ satisfies
\[
Q\cap F\neq\emptyset
\qquad\text{and}\qquad
Q-F\neq\emptyset.
\]
So the property is decidable by checking whether such a reachable subset exists.

\item \textbf{All strings have unique paths}: Construct a finite product machine
whose states record two simultaneous computations of $A$ on the same input,
together with a bit indicating whether the two computations have diverged. Then
there exists a string with two distinct complete paths iff a state with
``diverged = yes'' is reachable in this product machine. Hence uniqueness of
paths is decidable.
\end{enumerate}

%%------------------------------------------------------------
\item \textbf{Decidability for two DFAs $A_1$ and $A_2$.}

\begin{enumerate}
\item \textbf{Decidability of homomorphism existence}: A homomorphism $h: A_1\to A_2$ (of DFAs) is a state map preserving start state, transitions, and taking final states to final states. Check all functions $h: S_1\to S_2$: finitely many ($|S_2|^{|S_1|}$). For each, check the homomorphism conditions. Decidable (finite check).

\item \textbf{Decidability of isomorphism existence}: An isomorphism requires $h$ to be bijective (so $|S_1|=|S_2|$). Check all bijections: $(|S_1|)!$ many. For each, verify the isomorphism conditions. Decidable (finite check for finite $S_1$).

\item \textbf{Decidability of $>3$ isomorphisms}: Enumerate all bijections
$h:S_1\to S_2$. For each one, check whether it preserves the start state,
transitions, and final states. Count how many bijections pass these tests.
There are at most $(|S_1|)!$ candidates, so the procedure halts. Answer yes
exactly when the final count exceeds $3$. Hence the question is decidable. \qed
\end{enumerate}

%%------------------------------------------------------------
\item \textbf{Decidability: $R_1$ describes infinitely many strings (no DFA construction).}

\textbf{Algorithm} (structural induction on $R_1$):
- $R_1 = \pmb{\emptyset}$: finite (empty). No.
- $R_1 = \lambda$ or $R_1 = a$: finite. No.
- $R_1 = R \cup R'$: infinite iff $R$ or $R'$ is infinite.
- $R_1 = R\cdot R'$: infinite iff (R is infinite and $L(R')$ is nonempty) or ($L(R)$ is nonempty and $R'$ is infinite).
- $R_1 = R^*$: infinite iff $L(R)$ contains a nonempty string.

Each check is computable by structural recursion on $R_1$. Terminates since $R_1$ is finite. \qed

%%------------------------------------------------------------
\item \textbf{Decidability of Mealy machine equivalent to Moore machine.}

\textbf{Algorithm}:
\begin{enumerate}
\item Minimize the Mealy machine $M$ and the Moore machine $A$.
\item Convert one to the other's representation (Mealy to Moore or vice versa) using the standard transformation.
\item Check equivalence (Exercise 12.19 or 12.20). Decidable. \qed
\end{enumerate}

%%------------------------------------------------------------
\item \textbf{Decidability: $L(R_2)$ properly contains $L(R_1)$ ($L(R_1)\subset L(R_2)$).}

\textbf{Algorithm}:
\begin{enumerate}
\item Convert $R_1, R_2$ to DFAs.
\item Check $L(R_1)\subseteq L(R_2)$ (decidable, Exercise 12.4).
\item Check $L(R_1)\neq L(R_2)$ (decidable, Exercise 12.14 or 12.15).
\item $L(R_1)\subset L(R_2)$ iff both conditions hold. \qed
\end{enumerate}

%%------------------------------------------------------------
\item \textbf{Efficient method for $L(A)=\emptyset$ for NDFA $A$, without DFA conversion.}

\begin{enumerate}
\item \textbf{Method}: BFS/DFS directly on the NDFA $A$ (viewed as a graph). Start from $s_0$, compute the $\lambda$-closure, then follow transitions. $L(A)=\emptyset$ iff no final state is reachable from $s_0$ (via any sequence of transitions and $\lambda$-moves). Reachability on the finite graph of $A$ (without constructing the exponential subset DFA).

\item \textbf{Method for infinite $L(A)$}: First find all states reachable from
the initial $\lambda$-closure and all states from which some final state is
reachable. Keep only these useful states. Then compute the strongly connected
components of this useful subgraph. The language is infinite iff some useful
component contains a cycle with at least one non-$\lambda$ transition. Such a
cycle can be pumped to produce arbitrarily many different accepted strings; if
no such cycle exists, every accepting path has bounded input length. Hence the
question is decidable without converting to a DFA. \qed
\end{enumerate}

%%------------------------------------------------------------
\item \textbf{Decidability: $L(M)$ is a regular set for DPDA $M$.}

\textbf{Answer}: This question is decidable, but the proof is substantially
deeper than the earlier decidability arguments in the chapter. A standard
result on deterministic pushdown automata gives an algorithm that determines
whether the language of a given DPDA is regular. So the correct conclusion here
is simply that the problem is decidable.

\textit{Caution}: it is \emph{not} enough to ask whether this particular DPDA
happens to keep its stack bounded, since a DPDA can use unbounded stack space
and still recognize a regular language. \qed

%%------------------------------------------------------------
\item \textbf{Algorithm for path-finding in graphs (Theorem 12.9).}

\begin{enumerate}
\item \textbf{Appropriate algorithm}: BFS (breadth-first search) from the source node to find all reachable nodes and shortest paths. For directed graphs (as in the PDA product machine):
\begin{enumerate}[a.]
\item Initialize a queue with the start node; mark it visited.
\item At each step, dequeue a node, enqueue all unvisited neighbors, mark them visited.
\item Terminate when the queue is empty or the target is found.
\end{enumerate}
Time complexity: $O(|V| + |E|)$ where $V$ is vertices and $E$ is edges.

\item \textbf{More efficient recursive algorithm}: DFS with memoization.
\texttt{Reach}($v$, $t$): returns true if $t$ is reachable from $v$.
- If $v = t$: return true.
- If $v$ is marked: return false.
- Mark $v$; for each neighbor $u$ of $v$: if \texttt{Reach}($u$, $t$): return true.
- Return false.

This avoids reprocessing nodes and runs in $O(|V|+|E|)$ time. \qed
\end{enumerate}

%%------------------------------------------------------------
\item \textbf{\HSIG\ (halting-TM languages) closed under union, intersection, concatenation, reversal.}

\begin{enumerate}
\item \textbf{Union}: Given halting TMs $M_1$ ($L_1$) and $M_2$ ($L_2$), build $M$ for $L_1\cup L_2$: run $M_1$ and $M_2$ in parallel (interleaved simulation). Since both halt on all inputs, $M$ also halts. $M$ accepts iff $M_1$ or $M_2$ accepts. \qed

\item \textbf{Intersection}: Same parallel simulation; accept iff both accept. Both halt, so $M$ halts. \qed

\item \textbf{Concatenation}: $M$ for $L_1\cdot L_2$ on input $x$: enumerate all splits $x=uv$, run $M_1$ on $u$ and $M_2$ on $v$. Since both halt, each check halts. There are $|x|+1$ splits (finitely many), all decidable. Accept if any split works. \qed

\item \textbf{Reversal}: $M$ for $L_1^r$ on input $x$: compute $x^r$ (finitely computable), run $M_1$ on $x^r$. $M_1$ halts on $x^r$; so $M$ halts. Accept iff $M_1$ accepts $x^r$. \qed
\end{enumerate}

%%------------------------------------------------------------
\item \textbf{Halting TM $T'$ equivalent to $L_2(T)$ for halting TM $T$.}

$L = L_2(T)$: $T$ halts on all inputs and $L$ consists of strings for which $T$ halts with $\pmb{Y}$ on the tape.

Construct $T'$:
\begin{enumerate}[1.]
\item Simulate $T$ on input $w$.
\item $T$ halts (since it halts on all inputs by assumption).
\item If $T$ halted with $\pmb{Y}$ somewhere on the tape: erase the tape, write $w$ then $\pmb{Y}$ after $w$ on an otherwise blank tape; halt in the accepting state.
\item If $T$ halted with no $\pmb{Y}$: erase the tape, write $w$ then $\pmb{N}$ after $w$; halt in a rejecting state.
\end{enumerate}

$T'$ satisfies all three conditions: (1) halts on all input (since $T$ does and the post-processing is finite); (2) writes $\pmb{Y}$ after accepted words; (3) writes $\pmb{N}$ after rejected words. And $L(T') = L_2(T) = L$. \qed

%%------------------------------------------------------------
\item \textbf{Machine $M_{\mathfrak{M}}$ and application to Theorem 12.15.}

\begin{enumerate}
\item \textbf{Proof}: Assume there is a machine $M_{\mathfrak M}$ that decides
membership in $X$, and that every encoding in $X$ represents a halting Turing
machine. Define a halting Turing machine $D$ as follows on input $x$:

\begin{enumerate}
\item Use $M_{\mathfrak M}$ to test whether $x\in X$.
\item If $x\notin X$, reject.
\item If $x\in X$, let $T_x$ be the Turing machine encoded by $x$, simulate
$T_x$ on input $x$, and do the opposite: accept iff $T_x$ rejects $x$, and
reject iff $T_x$ accepts $x$.
\end{enumerate}

This simulation halts because every machine encoded by a member of $X$ is
halting. So $D$ is itself a halting Turing machine.

Now suppose $L(D)\in\mathfrak M$. Then there is some encoding $d\in X$ whose
machine $T_d$ recognizes $L(D)$. Consider what $D$ does on input $d$.

\begin{itemize}
\item If $d\in L(D)$, then by the definition of $D$, the machine $T_d$ rejects
$d$, contradicting $L(T_d)=L(D)$.
\item If $d\notin L(D)$, then by the definition of $D$, the machine $T_d$
accepts $d$, again contradicting $L(T_d)=L(D)$.
\end{itemize}

Thus $L(D)\notin\mathfrak M$. So there exists a halting Turing machine that
recognizes a language not in $\mathfrak M$. \qed

\item \textbf{Application to Theorem 12.15}: Let $X$ be the set of strings that
encode valid context-sensitive grammars by the scheme of the next exercise.
Membership in $X$ is decidable. Also, each member of $X$ represents a halting
Turing machine, because Theorem 12.9 gives an algorithm for deciding membership
in the language of a context-sensitive grammar. Therefore part (a) applies to
the class $\mathfrak M=\OSIG$ of context-sensitive languages, and yields a
halting Turing machine whose language is not context sensitive. That is exactly
Theorem 12.15.
\end{enumerate}

%%------------------------------------------------------------
\item \textbf{Encoding context-sensitive grammars.}

\begin{enumerate}
\item \textbf{Encoding scheme}: Assume 2 terminal symbols $\{a,b\}$ and arbitrarily many nonterminals $N_1, N_2, N_3, \ldots$

Encoding:
- Terminals: $a \mapsto \texttt{00}$, $b \mapsto \texttt{01}$.
- Nonterminal $N_i \mapsto \texttt{1}^i\texttt{0}$ (unary encoding with terminator).
- Productions: $\alpha\to\beta$ encoded as $\overline{\alpha}\texttt{\#\#}\overline{\beta}$ where $\overline{x}$ is the encoding of $x$ symbol by symbol.
- Grammar: list of productions separated by $\texttt{\#\#\#}$, preceded by the start symbol encoding.

This encoding handles an unbounded number of nonterminals via unary encoding and terminates on any finite grammar.

\item \textbf{Algorithm to decide valid CSG encoding}:
\begin{enumerate}
\item Parse the encoding to extract the alphabet, nonterminal list, start symbol, and productions.
\item Check that: (a) terminal symbols are from $\{a,b\}$; (b) nonterminals are distinct and properly encoded; (c) each production $\alpha\to\beta$ has $|\beta|\geq|\alpha|$ (context-sensitive condition), except possibly $Z\to\lambda$.
\item Verify the syntactic structure of the encoding (valid separators, well-formed strings).
\item The TM accepts valid encodings and rejects invalid ones. Terminates in polynomial time in the encoding length. \qed
\end{enumerate}
\end{enumerate}

%%------------------------------------------------------------
\item \textbf{Undecidability of $L(X)=\emptyset$.}

\begin{enumerate}
\item \textbf{Arbitrary Turing machines}: Undecidable. Given a Turing machine
$T$ and a word $w$, build a machine $M_{T,w}$ that ignores its own input,
simulates $T$ on $w$, and accepts iff that simulation halts. Then
\[
L(M_{T,w})=\emptyset
\]
iff $T$ does \emph{not} halt on $w$. So a decider for emptiness would solve the
halting problem.

\item \textbf{Halting Turing machines}: Still undecidable. Given an arbitrary
Turing machine $T$ and word $w$, build a halting Turing machine $H_{T,w}$ that
on input $x$ simulates $T$ on $w$ for exactly $|x|$ steps and accepts iff $T$
halts within that many steps. Then $H_{T,w}$ halts on every input. Also,
\[
L(H_{T,w})=\emptyset
\]
iff $T$ never halts on $w$. So deciding emptiness for halting Turing machines
would decide the halting problem.

\item \textbf{Context-sensitive grammars}: By Corollary 11.2, context-sensitive
grammars and LBAs define the same class of languages. So it is enough to prove
part (d): if emptiness were decidable for context-sensitive grammars, it would
also be decidable for LBAs by first converting the LBA to an equivalent grammar.

\item \textbf{Linear bounded automata}: Given $T$ and $w$, build an LBA
$B_{T,w}$ that on input $x$ of length $n$ simulates $T$ on $w$ for exactly $n$
steps and accepts iff $T$ halts within those $n$ steps. This can be done by a
linear bounded automaton because in $n$ steps the simulated machine can visit at
most $n+|w|$ tape cells, and $|w|$ is a fixed constant; the finitely many very
short inputs can be handled separately in the finite control.

Then
\[
L(B_{T,w})=\emptyset
\]
iff $T$ never halts on $w$. So emptiness for LBAs is undecidable, and therefore
emptiness for context-sensitive grammars is undecidable as well. \qed
\end{enumerate}

%%------------------------------------------------------------
\item \textbf{Undecidability of $L(X)=\Sigma^*$.}

\begin{enumerate}
\item \textbf{Arbitrary TMs}: Undecidable. Given $T$ and $w$, build a machine
$M'_{T,w}$ that ignores its input, simulates $T$ on $w$, and accepts iff $T$
halts. Then
\[
L(M'_{T,w})=\Sigma^*
\]
iff $T$ halts on $w$. So universality for arbitrary Turing machines is
undecidable.
\item \textbf{Halting TMs}: Undecidable. Given $T$ and $w$, define a halting
Turing machine $U_{T,w}$ that on input $x$ simulates $T$ on $w$ for exactly
$|x|$ steps and accepts iff $T$ has \emph{not} halted within that many steps.
Then $U_{T,w}$ halts on all inputs, and
\[
L(U_{T,w})=\Sigma^*
\]
iff $T$ never halts on $w$. So universality for halting Turing machines is
undecidable.
\item \textbf{Context-sensitive grammars}: By Corollary 11.2, it is enough to
prove the corresponding statement for LBAs.
\item \textbf{Linear bounded automata}: Given $T$ and $w$, build an LBA
$C_{T,w}$ that on input $x$ of length $n$ simulates $T$ on $w$ for exactly $n$
steps and accepts iff $T$ has \emph{not} halted within those $n$ steps. The
simulation uses only linear space: the machine keeps $x$ on one track, uses a
counter of size $O(\log n)$ to count from $0$ up to $n$, and stores on the
remaining track space a length-$O(n)$ window describing the first $n$ tape cells
that $T$ could have visited in its first $n$ steps. Since a machine running for
only $n$ steps cannot move farther than $n$ cells from its start position, this
bounded work area is sufficient. Hence $C_{T,w}$ is indeed an LBA.
Now
\[
L(C_{T,w})=\Sigma^*
\]
iff $T$ never halts on $w$. Hence universality is undecidable for LBAs, and
therefore also for context-sensitive grammars.
\item \textbf{Context-free grammars}: Undecidable. This is a standard
undecidability result for context-free grammars.
\item \textbf{Pushdown automata}: Undecidable, since Theorem 10.3 gives an
effective conversion between PDAs and context-free grammars. \qed
\end{enumerate}

%%------------------------------------------------------------
\item \textbf{Set $E$ of encodings of TMs halting on $\lambda$.}

\begin{enumerate}
\item \textbf{$E\in\TSIG$}: Yes. The TM that recognizes $E$ works as follows:
on input $\langle T\rangle$, simulate $T$ on $\lambda$. If $T$ halts, accept. If
$T$ does not halt, keep simulating forever. Therefore $E$ is Turing acceptable.

\item \textbf{$E\in\HSIG$}: No. If $E\in\HSIG$, then $E$ would be decidable. But
$E$ is exactly the halting problem on input $\lambda$, which is undecidable. So
$E\notin\HSIG$.

\item \textbf{$E\in\OSIG$}: No. By Corollary 12.3, every context-sensitive
language belongs to $\HSIG$. Since $E\notin\HSIG$, it certainly cannot belong to
$\OSIG$.
\end{enumerate}

%%------------------------------------------------------------
\item \textbf{Set $N$ of encodings of TMs that do NOT halt on $\lambda$.}

\begin{enumerate}
\item \textbf{$N\in\TSIG$}: No. Here
\[
N=\overline{E}.
\]
If both $E$ and $\overline{E}$ belonged to $\TSIG$, then we could dovetail the
two semidecision procedures and decide $E$, which would put $E$ in $\HSIG$.
But Exercise 12.41(b) showed that $E\notin\HSIG$. Therefore $N\notin\TSIG$.

\item \textbf{$N\in\HSIG$}: No. If $N$ were in $\HSIG$, then both $E$ and
$N=\overline{E}$ would be recursive languages. A recursive language and its
complement are both decidable, so $E$ would be decidable. But the non-halting
problem is undecidable. Therefore $N\notin\HSIG$.

\item \textbf{$N\in\OSIG$}: No. Again, $\OSIG\subseteq\HSIG$, and we have just
seen that $N\notin\HSIG$. Hence $N\notin\OSIG$. \qed
\end{enumerate}

%%------------------------------------------------------------
\item \textbf{Decidability questions about equivalence classes of $R_{L(A)}$.}

\begin{enumerate}
\item \textbf{At least one equivalence class is finite}: Minimize $A$ to obtain
the connected DFA
\[
M=\langle \Sigma,S,s_0,\delta,F\rangle.
\]
Its states correspond to the classes of $R_{L(A)}$. For a state $q\in S$, the
corresponding class is
\[
I(M,q)=\{x\in\Sigma^*\mid \DBAR(s_0,x)=q\}.
\]
This set is infinite iff some directed cycle reachable from $s_0$ can also
reach $q$: traversing such a cycle any number of times gives infinitely many
different strings in $I(M,q)$. Conversely, if no reachable cycle can reach
$q$, then only finitely many paths lead from $s_0$ to $q$, so $I(M,q)$ is
finite.

Therefore at least one class is finite iff there is some state $q$ not
reachable from any reachable cycle. This is decidable by graph search.

\item \textbf{All equivalence classes are finite}: By the same criterion, every
class is finite iff no reachable cycle exists at all in the minimal DFA. So it
suffices to minimize $A$ and test whether the reachable graph is acyclic. This
is decidable. \qed
\end{enumerate}

%%------------------------------------------------------------
\item \textbf{Decide $L(G_1)\subset L(G_2)$ for right-linear grammars $G_1$, $G_2$.}

\textbf{Algorithm}:
\begin{enumerate}
\item Convert $G_1, G_2$ to DFAs $A_1, A_2$ (via Lemma 8.1).
\item Check $L(A_1)\subseteq L(A_2)$ (decidable, Exercise 12.4).
\item Check $L(A_1)\neq L(A_2)$ (decidable, Exercise 12.14).
\item $L(G_1)\subset L(G_2)$ (proper subset) iff $L(A_1)\subseteq L(A_2)$ and $L(A_1)\neq L(A_2)$. \qed
\end{enumerate}

\end{enumerate}
